\section{General background}
\label{sec:background}

\subsection{Preliminary notational and geometric definitions}

The main geometric objects used along the present work are shown in Figure \ref{fig:geometry}.
We consider a topocentric Cartesian coordinate system with axes 
$x$, $y$, and $z$ pointing to north, east and down, respectively. 
We also assume to have a set of 3D gravity or magnetic sources with top at the vertical coordinate $z^{t}_{0}$, 
located in the subsurface over $\overline{\mathcal{V}} = \mathcal{V} \, \cup \, \partial \mathcal{V}$, where 
$\mathcal{V} \subset \mathbb{R}^{3}$ is a \textit{regular region} with boundary
$\partial \mathcal{V}$. 
Some important properties of $\mathcal{V}$ in this case is that it is an open connected set of points 
and $\partial \mathcal{V}$ is a closed and compact surface free of double points \citep[][ p. 22]{freeden_gerhards2012}.
Generic points inside the sources are represented by $\mathbf{r}' = (x', y', z')^{\top} \in \overline{\mathcal{V}}$.

We also consider to have potential field observations at generic points 
$\mathbf{r} = (x, y, z)^{\top} \notin \overline{\mathcal{V}}$ or at a finite set of points 
$\mathbf{r}_{i} = (x_{i}, y_{i}, z_{i})^{\top} \notin \overline{\mathcal{V}}$, $i \in \{ 1, \dots, N \}$, 
located above the sources, i.e., with vertical coordinates satisfying $\forall z < z^{t}_{0}$ and 
$\forall z_{i} < z^{t}_{0}$. These points are conveniently called \textit{observation points}.

Specific horizontal surfaces are defined by $\mathcal{S}$, with a particular superscript symbol. 
The horizontal surface $\mathcal{S}''$, for example, is defined at a vertical coordinate $z''_{0}$ below the 
observation points, so that $\forall z < z''_{0}$ and $\forall z_{i} < z''_{0}$. 
The subscript ``$0$'' indicates that the corresponding vertical coordinate is constant. 
Points and elements of area on a given horizontal surface are defined by using the corresponding superscript symbol. 
For example, we define the element of area $d\mathcal{S}'' = dx''dy''$, generic points $\mathbf{r}'' = (x'', y'', z''_{0})^{\top} \in \mathcal{S}''$ 
and also a finite set of points $\mathbf{r}''_{j} = (x''_{j}, y''_{j}, z''_{0})^{\top} \in \mathcal{S}''$,
$j \in \{ 1, \dots, M \}$.

Finally, all Fourier transforms computed here are defined on horizontal surfaces, with respect to the 
spatial coordinates along the horizontal directions. For example, $\mathcal{F} [ U(\mathbf{r}^{\dagger}) ]$
indicates the two-dimensional Fourier transform of a scalar function $U(\mathbf{r}^{\dagger})$ with respect to
the spatial coordinates $x$ and $y$, on a horizontal surface $\mathcal{S}^{\dagger}$ formed by points  
$\mathbf{r}^{\dagger} = (x^{\dagger}, y^{\dagger}, z^{\dagger}_{0})^{\top}$, i.e. by $\mathbf{r}^{\dagger} \in \mathcal{S}^{\dagger}$.
Besides, the dependence of $\mathcal{F} [ U(\mathbf{r}^{\dagger}) ]$ on the wavenumbers $k_{x}$ and $k_{y}$ associated with 
$x$ and $y$ directions, respectively, is conveniently omitted to simplify the notation.

\subsection{Potential fields due to 3D sources}
\label{sec:potential_fields_3D_sources}

We define the potential field $d^{o}(\mathbf{r})$ produced by the 3D sources at observation points 
$\mathbf{r}$ (Figure \ref{fig:geometry}) by
\begin{equation}
	d^{o}(\mathbf{r}) = D_{d} \, U(\mathbf{r}) \: , 
	\quad \mathbf{r} \notin \overline{\mathcal{V}} \: ,
	\label{eq:observed-data-space}
\end{equation}
where $D_{d}$ is a space-domain differential operator acting on the potential
\begin{equation}
	U(\mathbf{r}) = \kappa \iiint\limits_{\mathcal{V}}
	\sigma(\mathbf{r}') \: \frac{1}{\| \mathbf{r} - \mathbf{r}'\|}
	\: d\mathcal{V}' \: , 
	\quad \mathbf{r} \notin \overline{\mathcal{V}} \: ,
	\label{eq:potential-U-space}
\end{equation}
$d\mathcal{V}' = dx'dy'dz'$ represents a volume element;
$\kappa$ is a proportionality constant;
$\tfrac{1}{\| \mathbf{r} - \mathbf{r}' \|}$ is the inverse distance function
between the observation point $\mathbf{r}$ and a point $\mathbf{r}' \in \mathcal{V}$
and $\sigma(\mathbf{r}')$ is a scalar function defined at points 
$\mathbf{r}' = (x', y', z')^{\top} \in \overline{\mathcal{V}}$ 
as the density (contrast) or total-magnetization intensity within the sources.
We implicitly assume that $\sigma(\mathbf{r}) = 0$ outside $\overline{\mathcal{V}}$.
To insure that the partial derivatives of first and second orders of $U$ 
exist at all points in the $\mathbb{R}^{3}$ space, including those within the sources,
we impose that $\sigma(\mathbf{r}')$ is a piecewise continuous, integrable and bounded function
that satisfies some Hölder condition \citep[][ p. 152]{kellogg1967} at all points
$\mathbf{r}' \in \overline{\mathcal{V}}$.

The Fourier transform of the potential field $d^{o}$ (equation \ref{eq:observed-data-space}) 
on a horizontal surface above the sources is given by \citep[e.g.,][]{bhattacharyya1967,gunn1975}:
\begin{equation}
	\mathcal{F} [ d^{o}(\mathbf{r}^{\dagger}) ] =
	\alpha_{d} \, \mathcal{F} [ U(\mathbf{r}^{\dagger}) ] \: ,
	\quad \mathbf{r}^{\dagger} \in \mathcal{S}^{\dagger} \: ,
	\quad z^{\dagger}_{0} < z^{t}_{0} \: ,
	\label{eq:observed-data-Fourier}
\end{equation}
where 
\begin{equation}
	\mathcal{F} [ U(\mathbf{r}^{\dagger}) ] = 2 \pi \kappa 
	\int\limits_{z^{t}_{0}}^{\infty}
	\mathcal{F} [ \sigma(\mathbf{r}') ] \:
	\frac{e^{- k \, (z' - z^{\dagger}_{0})}}{k}
	\: dz' \: , 
	\quad \mathbf{r}^{\dagger} \in \mathcal{S}^{\dagger} \: ,
	\quad z^{\dagger}_{0} < z^{t}_{0} \: ,
	\label{eq:potential-U-Fourier}
\end{equation}
is the Fourier transform of the Newtonian potential 
$U$ (equation \ref{eq:potential-U-space}) on $\mathcal{S}^{\dagger}$, with $k = \sqrt{k_{x}^{2} + k_{y}^{2}}$, and
\begin{equation}
	\mathcal{F} [ \sigma(\mathbf{r}') ] = 
	\int\limits_{-\infty}^{\infty}\int\limits_{-\infty}^{\infty}
	\sigma(\mathbf{r}') \: e^{-\mathtt{i} \, (k_{x} \, x' + k_{y} \, y')}
	\: dx'dy'
	\label{eq:sigma-Fourier}
\end{equation}
is the Fourier transform of $\sigma(\mathbf{r}')$ 
(equation \ref{eq:potential-U-space}) on the horizontal plane at the vertical
coordinate $z'$ and $\mathtt{i}$ represents the imaginary unit.
In equation \ref{eq:observed-data-Fourier}, $\alpha_{d}$ is a function of the 
wavenumbers $k_{x}$ and $k_{y}$ and represents the differential operator 
$D_{d}$ (equation \ref{eq:observed-data-space}) in the Fourier domain.

The differential operators $D_{d}$ and $\alpha_{d}$ associated with the potential fields commonly found 
in geophysical applications (equations \ref{eq:observed-data-space} and \ref{eq:observed-data-Fourier}) 
can be defined as follows:
\begin{equation}
	D(\ell, \hat{\mathbf{u}}) = 
	\left( \hat{\mathbf{u}}^{\top} \nabla \right)^{\ell} \Leftrightarrow
	\alpha(\ell, \hat{\mathbf{u}}) = \left( \hat{\mathbf{u}}^{\top} \mathbf{k} \right)^{\ell} \: ,
	\quad \ell \in \{ 0, 1, 2 \} \: ,
	\label{eq:diff-operators-generic}
\end{equation}
where $\hat{\mathbf{u}}$ is a constant unit vector,
$\nabla$ is the vector differential operator nabla \citep[e.g.,][ p. 179]{kaplan2002} and 
$\mathbf{k} = ( \mathtt{i} \, k_{x} \, , \, \mathtt{i} \, k_{y} \, , \, k)^{\top}$
\citep{pedersen1978}.
The differential operator $D$ (equation \ref{eq:diff-operators-generic})
has a symbolic form based on that used by \citet[][ p. 179]{kaplan2002}, for example, where the suggested
dot product and exponentiation actually leads to a differentiation.

By combining successive applications of $D(\ell, \hat{\mathbf{u}})$ (equation \ref{eq:diff-operators-generic})
with properly chosen $\ell$ and $\hat{\mathbf{u}}$, we may define differential operators that 
result in all potential fields $d^{o}$ commonly used in Geophysics.
Note that $D(0, \hat{\mathbf{u}}) = \alpha(0, \hat{\mathbf{u}}) = 1$ produces no differentiation.
It can be verified, for example, that using $\ell = 1$ and a constant unit vector $\hat{\mathbf{u}}$
in equation \ref{eq:diff-operators-generic} leads to the following directional derivative 
of first order
\begin{equation}
	D(1, \hat{\mathbf{u}}) \, U(\mathbf{r}) = \nabla U(\mathbf{r})^{\top} \hat{\mathbf{u}}
	\label{eq:potential-U-oblique-1st-deriv-space}
\end{equation}
where
\begin{equation}
	\nabla U(\mathbf{r}) = \left( 
	\tfrac{\partial}{\partial x} U(\mathbf{r}),
	\tfrac{\partial}{\partial y} U(\mathbf{r}),
	\tfrac{\partial}{\partial z} U(\mathbf{r}) \right)^{\top}
	\label{eq:potential-U-grad-space}
\end{equation}
is the gradient of $U(\mathbf{r})$ (equation \ref{eq:potential-U-space}), with elements
given by the partial derivatives of first order
\begin{equation}
	\tfrac{\partial}{\partial u} U(\mathbf{r}) = \kappa \iiint\limits_{\mathcal{V}}
	\sigma(\mathbf{r}') \: 
	\tfrac{\partial}{\partial u} \left( \frac{1}{\| \mathbf{r} - \mathbf{r}'\|} \right)
	\: d\mathcal{V}' \: , 
	\quad \mathbf{r} \notin \overline{\mathcal{V}} \: , \quad u \in \{ x, y, z \} \: .
	\label{eq:potential-U-1st-deriv-space}
\end{equation}
The first-order directional derivative (equation \ref{eq:potential-U-oblique-1st-deriv-space})
coincides with the partial derivatives of first order $\tfrac{\partial}{\partial u} U$ (equation \ref{eq:potential-U-1st-deriv-space})
for the particular case in which the unit vector $\hat{\mathbf{u}}$ coincides with one of the canonical base vectors
\begin{equation}
	\hat{\mathbf{x}} = (1,0,0)^{\top} , \quad 
	\hat{\mathbf{y}} = (0,1,0)^{\top} , \quad
	\hat{\mathbf{z}} = (0,0,1)^{\top} \: .
	\label{eq:canonical-base-vector}
\end{equation}
For example, $D(1, \hat{\mathbf{z}}) \, U(\mathbf{r}) = \tfrac{\partial}{\partial z} U(\mathbf{r})$.

It can also be verified from equation \ref{eq:diff-operators-generic} that 
\begin{equation}
	D(1, \hat{\mathbf{u}}) \, D(1, \hat{\mathbf{v}}) \, U(\mathbf{r}) = 
	\hat{\mathbf{u}}^{\top} \mathbf{H}_{U}(\mathbf{r}) \, \hat{\mathbf{v}} \: ,
	\label{eq:potential-U-oblique-2nd-deriv-space}
\end{equation}
represents a second-order directional derivative along constant unit vectors $\hat{\mathbf{u}}$ and
$\hat{\mathbf{v}}$, where 
\begin{equation}
	\mathbf{H}_{U}(\mathbf{r}) = 
	\begin{bmatrix}
		\tfrac{\partial^{2}}{\partial x^{2}} U(\mathbf{r}) & 
		\tfrac{\partial^{2}}{\partial x \, \partial y} U(\mathbf{r}) & 
		\tfrac{\partial^{2}}{\partial x \, \partial z} U(\mathbf{r}) \\
		\tfrac{\partial^{2}}{\partial x \, \partial y} U(\mathbf{r}) & 
		\tfrac{\partial^{2}}{\partial y^{2}} U(\mathbf{r}) & 
		\tfrac{\partial^{2}}{\partial y \, \partial z} U(\mathbf{r}) \\
		\tfrac{\partial^{2}}{\partial x \, \partial z} U(\mathbf{r}) & 
		\tfrac{\partial^{2}}{\partial y \, \partial z} U(\mathbf{r}) & 
		\tfrac{\partial^{2}}{\partial z^{2}} U(\mathbf{r})
	\end{bmatrix}
	\label{eq:potential-U-Hessian}
\end{equation}
is the Hessian matrix \citep[e.g.,][ p. 158]{kaplan2002} of $U$ (equation \ref{eq:potential-U-space}),
with elements given by the partial derivatives of second order
\begin{equation}
	\tfrac{\partial^{2}}{\partial u \, \partial v} U(\mathbf{r}) = \kappa \iiint\limits_{\mathcal{V}}
	\sigma(\mathbf{r}') \: 
	\tfrac{\partial^{2}}{\partial u \, \partial v} \left( \frac{1}{\| \mathbf{r} - \mathbf{r}'\|} \right)
	\: d\mathcal{V}' \: , \quad \mathbf{r} \notin \overline{\mathcal{V}} \: , \quad u, v \in \{ x, y, z \} \: .
	\label{eq:potential-U-2nd-deriv-space}
\end{equation} 
It can be easily verified that 
the second-order directional derivative $D(1, \hat{\mathbf{u}}) \, D(1, \hat{\mathbf{v}}) \, U(\mathbf{r})$ 
(equation \ref{eq:potential-U-oblique-2nd-deriv-space})
coincides with the partial derivatives of second order $\tfrac{\partial^{2}}{\partial u \, \partial v} U(\mathbf{r})$ 
(equation \ref{eq:potential-U-2nd-deriv-space}) for the particular cases in which the unit vectors 
$\hat{\mathbf{u}}$ and $\hat{\mathbf{v}}$ are given by the canonical base vectors (equation \ref{eq:canonical-base-vector}).
For example, $D(1, \hat{\mathbf{x}}) \, D(1, \hat{\mathbf{y}}) \, U(\mathbf{r}) = \tfrac{\partial^{2}}{\partial x \, \partial y} U(\mathbf{r})$.

\subsection{Some relevant potential fields for geophysical applications}
\label{subsec:relevant-potential-fields}

In gravity exploration methods, the vertical derivative $\tfrac{\partial}{\partial z} U$
(equation \ref{eq:potential-U-1st-deriv-space}) does not define the gravity anomaly, but the \textit{gravity disturbance} data 
\cite[e.g.,][ p. 93]{hofmann-wellenhof-moritz2005}.
The debate around the differences between gravity anomaly and gravity disturbance is beyond the scope of the 
present work and has been carried about from a more geodetic than a geophysical point of view 
\citep[e.g.,][]{li2001, fairhead2003, hackney-featherstone2003, hinze2005, vajda-etal2006, vajda-etal2007, vajda-etal2008}. 
Maybe, that is why it is still unfamiliar to a considerable group of exploration geophysicists.
Another important quantity for gravity exploration method is the Hessian matrix $\mathbf{H}_{U}$
(equation \ref{eq:potential-U-Hessian}), which defines the
(gravitational) \textit{gradient tensor} \citep[e.g.,][]{pedersen_rasmussen1990}.

In magnetic exploration methods, it is well-established that the second-order directional derivative 
$\hat{\mathbf{u}}^{\top} \mathbf{H}_{U} \, \hat{\mathbf{v}}$ (equation \ref{eq:potential-U-oblique-2nd-deriv-space})
is commonly used to approximate the \textit{total-field anomaly} produced by 3D sources \cite[e.g., ][ p. 179]{blakely1996},
where $\hat{\mathbf{u}}$ defines the constant direction of the main geomagnetic
field at the study area and $\hat{\mathbf{v}}$ the constant total-magnetization direction of the sources.
Because of that, we use the term \textit{approximated total-field anomaly} to denote this 
potential field throughout the manuscript.
Note that, in this case, $\hat{\mathbf{u}}^{\top} \mathbf{H}_{U} \, \hat{\mathbf{v}}$ 
(equation \ref{eq:potential-U-oblique-2nd-deriv-space}) can be used to define (i) the  
\textit{Cartesian components of the magnetic induction field} produced due to the 3D
sources by setting the unit vector $\hat{\mathbf{u}}$ equal to the canonical base vectors $\hat{\mathbf{x}}$,
$\hat{\mathbf{y}}$ and $\hat{\mathbf{z}}$ (equation \ref{eq:canonical-base-vector}) or (ii)
the \textit{RTP anomaly} by setting $\hat{\mathbf{u}} = \hat{\mathbf{v}} = \pm\hat{\mathbf{z}}$.
This special anomaly represents the vertical component of 
the magnetic induction field that would be produced by the 3D sources if they had a 
purely induced and vertical magnetization. Differently from generic total-field anomalies 
\cite[e.g.,][ p. 178]{blakely1996}, which have a ``skewed'' pattern 
oscillating from negative to positive values and do not have their maxima placed right 
above the magnetic sources, the RTP anomaly is easier to interpret 
because it is predominantly positive right above the sources and
decays to zero towards their boundaries. \citet{baranov1957} was the first
to propose a method, called \textit{reduction to the pole} (RTP), to transform generic
total-field anomalies into this special anomaly.

\subsection{Classical Dirichlet's and Neumann's problems on a plane}
\label{subsec:dirichlet_neumann}

At all points $\mathbf{r} \notin \overline{\mathcal{V}}$ 
located outside the sources, $U$ (equation
\ref{eq:potential-U-space}) has continuous partial derivatives of all orders, 
satisfies the Laplace's equation and is regular at infinity \citep[][ p. 122--124, 217]{kellogg1967}.
Because of these properties, we have that
\begin{equation}
	U(\mathbf{r}) = 
	\int\limits_{-\infty}^{\infty} \int\limits_{-\infty}^{\infty}
	D(\ell, \hat{\mathbf{z}}) U(\mathbf{r}'') \:\: 
	D(1 - \ell, \hat{\mathbf{z}}) \Psi(\mathbf{r}, \mathbf{r}'')
	\: d\mathcal{S}'' \: , \quad \forall z < z''_{0} < z^{t}_{0} \: ,
	\quad \ell \in \{0, 1\} \: ,
	\label{eq:U-Dirichlet-Neumann}
\end{equation}
where $D(\ell, \hat{\mathbf{u}})$ is the space-domain differential operator 
defined by equation \ref{eq:diff-operators-generic} in terms of the canonical base vector 
$\hat{\mathbf{z}}$ (equation \ref{eq:canonical-base-vector}) and 
\begin{equation}
	\Psi(\mathbf{r}, \mathbf{r}'') = 
	\frac{1}{2\pi} \: 
	\frac{1}{\| \mathbf{r} - \mathbf{r}'' \|} \: , 
	\quad \mathbf{r} \ne \mathbf{r}'' \: ,
	\label{eq:Psi}
\end{equation}
represents the scaled inverse distance function between $\mathbf{r}$ and $\mathbf{r}''$. 
Equation \ref{eq:U-Dirichlet-Neumann} is commonly deduced from 
the \textit{Green's third identity} \citep[][ p. 219]{kellogg1967}, according to
the general outline presented by \cite{roy1962}, for example.

For $\ell = 0$, equation \ref{eq:U-Dirichlet-Neumann} states that $U$ can be defined
in terms of its values $D(0, \hat{\mathbf{u}}) U(\mathbf{r}'') = U(\mathbf{r}'')$ on a horizontal surface $\mathcal{S}''$ 
above the sources and 
represents the solution of the classical \textit{Dirichlet’s problem} or the 
\textit{first boundary value problem of potential theory} \cite[][ p. 236]{kellogg1967} on a plane, where 
$D(1, \hat{\mathbf{z}}) \Psi(\mathbf{r}, \mathbf{r}'') = \tfrac{\partial}{\partial z} \Psi(\mathbf{r}, \mathbf{r}'')$ 
represents the Green's function.
From the physical point of view, equation \ref{eq:U-Dirichlet-Neumann} with $\ell = 0$ defines a continuous 
layer of monopoles (or vertically-magnetized dipoles) that is located on $\mathcal{S}''$, has 
physical-property distribution $U(\mathbf{r}'')$ and exactly reproduces the 
potential field $U(\mathbf{r})$ at any point $\mathbf{r}$ above $\mathcal{S}''$.

On the other hand, equation \ref{eq:U-Dirichlet-Neumann} for $\ell = 1$ shows that $U$ can also
be defined in terms of its vertical derivative 
$D(1, \hat{\mathbf{z}}) U(\mathbf{r}'') = \tfrac{\partial}{\partial z} U(\mathbf{r}'')$ 
a horizontal surface $\mathcal{S}''$ above the sources and represents the solution of the classical \textit{Neumann’s problem} or the 
\textit{second boundary value problem of potential theory} \cite[][ p. 246]{kellogg1967} on a plane, 
with the Green's function given by $D(0, \hat{\mathbf{u}}) \Psi(\mathbf{r}, \mathbf{r}'') = \Psi(\mathbf{r}, \mathbf{r}'')$.
Equation \ref{eq:U-Dirichlet-Neumann} with $\ell = 1$ can be alternatively interpreted as a planar layer of
vertically-magnetized dipoles (or monopoles) that is located on $\mathcal{S}''$, has 
physical-property distribution $\tfrac{\partial}{\partial z} U(\mathbf{r}'')$ and exactly reproduces 
$U(\mathbf{r})$ at any point $\mathbf{r}$ above $\mathcal{S}''$.

By applying the Fourier-convolution theorem \cite[][ p. 268]{blakely1996} to equation \ref{eq:U-Dirichlet-Neumann} 
on a horizontal surface $\mathcal{S}^{\dagger}$ above the sources, we find that
\begin{equation}
	\mathcal{F} [ U(\mathbf{r}^{\dagger}) ] = 
	k \: \mathcal{F} [ U(\mathbf{r}'') ] \:
	\mathcal{F} [ \Psi(\mathbf{r}^{\dagger}, \mathbf{r}'') ] \: ,
	\quad z^{\dagger}_{0} < z''_{0} < z^{t}_{0} \: , \quad k \ne 0 \: ,
	\label{eq:U-Dirichlet-Neumann-Fourier}
\end{equation}
where 
\begin{equation}
	\mathcal{F} [ \Psi(\mathbf{r}^{\dagger}, \mathbf{r}'') ] 
	= \frac{e^{- k \, (z''_{0} - z^{\dagger}_{0})}}{k}
	\: , \quad z^{\dagger}_{0} < z''_{0} \: , \quad k \ne 0 \: ,
	\label{eq:Psi-Fourier}
\end{equation}
is defined according to \citet[][ p. 275]{blakely1996}, with $k = \sqrt{k_{x}^{2} + k_{y}^{2}}$.
We get from equation \ref{eq:diff-operators-generic} that the term $k$ in equation \ref{eq:U-Dirichlet-Neumann-Fourier} represents,
in the Fourier domain, the combined effect of the space-domain differential operators $D(\ell, \hat{\mathbf{z}})$ and 
$D(1-\ell, \hat{\mathbf{z}})$.

\subsection{Potential-field transformations}
\label{subsec:potential_field_transformations}

Several geophysical applications use the observed potential field 
$d^{o}$ (equation \ref{eq:observed-data-space}) to estimate another potential field 
defined in the space domain by
\begin{equation}
	f(\mathbf{r}) = D_{f} \, U(\mathbf{r}) \: , 
	\quad \mathbf{r} \notin \overline{\mathcal{V}} \: ,
	\label{eq:transformed-field-space}
\end{equation}
or, in the Fourier domain, 
\begin{equation}
	\mathcal{F} [ f(\mathbf{r}^{\dagger}) ] =
	\alpha_{f} \, \mathcal{F} [ U(\mathbf{r}^{\dagger}) ] \: ,
	\quad \mathbf{r}^{\dagger} \in \mathcal{S}^{\dagger} \: ,
	\quad z^{\dagger}_{0} < z^{t}_{0} \: ,
	\label{eq:transformed-field-Fourier}
\end{equation}
where $D_{f}$ is a space-domain differential operator acting on $U$ (equation \ref{eq:potential-U-space}) and
$\alpha_{f}$ its corresponding representation in the Fourier domain.
The differential operators $D_{f}$ and $\alpha_{f}$ associated with the transformed potential field 
$f$ (equations \ref{eq:transformed-field-space} and \ref{eq:transformed-field-Fourier}) are commonly
defined in terms of $D$ and $\alpha$ (equation \ref{eq:diff-operators-generic}). 
Estimating $f$ (equations \ref{eq:transformed-field-space} and \ref{eq:transformed-field-Fourier}) from 
the observed potential field $d^{o}$ (equation \ref{eq:observed-data-space}) is an operation 
commonly known in geophysics as \textit{potential-field transformation}.

A common approach to compute potential-field transformations \cite[e.g.,][]{gunn1975} 
consists in using the definition of observed data in the Fourier domain (equation \ref{eq:observed-data-Fourier}) 
and manipulating the Fourier domain solution of the Dirichlet's and Neumann's problems 
(equation \ref{eq:U-Dirichlet-Neumann-Fourier}) to obtain
\begin{equation}
	\mathcal{F} [ f(\mathbf{r}^{\dagger}) ] =
	\left( \tfrac{\alpha_{f}}{\alpha_{d}} \right) \, \mathcal{F} \left[ d^{o}(\mathbf{r}^{\dagger}) \right] \: , 
	\quad \mathbf{r}^{\dagger} \in \mathcal{S}^{\dagger} \: ,
	\quad z^{\dagger}_{0} < z^{t}_{0} \: .
	\label{eq:field-transformation-Fourier}
\end{equation}
Equation \ref{eq:field-transformation-Fourier} defines a single-step potential-field transformation 
in terms of the observed potential field $d^{o}(\mathbf{r}^{\dagger})$ on a planar surface $\mathcal{S}^{\dagger}$ 
located above the sources, at $z^{\dagger}_{0}$. 
The transformed field $f$ in the space domain is obtained by simply computing the inverse Fourier 
transform of \ref{eq:field-transformation-Fourier}.
For the sake of simplicity, we have opted for defining the fields $d^{o}$ and $f$ at the same points in equation 
\ref{eq:field-transformation-Fourier}.
However, it is important to stress that, depending on the 
geophysical application, $f$ may be estimated at points different from those at which the observations are 
placed.

The \textit{equivalent-layer technique} offers an alternative approach that is generally defined in the 
space domain to compute potential-field transformations. What follows aims at discussing the well-established 
empirical approaches used in the equivalent-layer technique, as well as discussing the theoretical background 
supporting them.